\emph{Odd-sector map of the non-centred Gaussian rank-one edge charge.} Within the Gaussian state model of the second paragraph above, relaxing its leading-order convention to a Gaussian \(\operatorname{Im}\alpha\) of mean \(\bar a\) and variance \(V\) independent of the slot radii, and optionally displacing the exponential slots by shifts \(c_i\ge0\), all averaged charge moments close exactly at every order: for every integer \(m\ge1\), \(\langle\operatorname{Tr}Q^{m}\rangle=(-1)^{m}2^{m+1}\sum_{k\le m/2}\binom{m}{2k}\sum_{j=0}^{k}\binom{k}{j}A_{m-2k+2j}H_{k-j}\), where \(A_n=\mathbb{E}\bigl((\operatorname{Im}\alpha)^{n}\bigr)\) are non-centred normal moments and \(H_n=\langle r^{n}\rangle\), verified symbolically against the exact trace polynomials and against direct double integrals at rational data. For odd \(m=2p+1\) this is an odd polynomial in \(\bar a\) with leading term \(\langle\operatorname{Tr}Q\rangle=-4\bar a\), so in the centred model every odd averaged moment vanishes identically for every variance and every slot law independent of \(\operatorname{Im}\alpha\): the odd-sector vanishing asserted in the second paragraph above is a symmetry statement, not a consequence of neglecting \(\operatorname{Im}\alpha\). The slot shifts enter only through the moment kernel \(\prod_i(1-s^{2}g_i t)^{-1}\exp\bigl(c_i s^{2}g_i t/(1-s^{2}g_i t)\bigr)\) for \(H_n=n!\) times the coefficient of \(t^{n}\), equivalently \(H_n=\sum_{\pi}\prod_{B\in\pi}\kappa_{|B|}\) over set partitions of \(\{1,\dots,n\}\) with cumulants \(\kappa_k=(k-1)!\,s^{2k}\sum_i g_i^{k}(1+k c_i)\), the real-graph presentation calculus of Kuriki and Numata specialised to this rank-one law; the shifts never by themselves generate odd charge moments while the slots stay independent of \(\operatorname{Im}\alpha\), and at \(\bar a=0\), \(V=0\) the even half of the map reproduces the closed moment map of the second paragraph above exactly. The mean, the variance and the slot shifts are inputs of the named state class, they do not fix the state that sets \(W\) and \(g\), leave the missing steps (i)-(iv) of the Missing step paragraph unchanged, and they close nothing in Problem~G.
